\documentclass[11pt,letterpaper]{article}

\usepackage[utf8]{inputenc}
\usepackage[spanish]{babel}
\usepackage[margin=2.5cm]{geometry}
\usepackage{amsmath,amssymb,amsthm}
\usepackage{enumitem}
\usepackage{titlesec}
\usepackage{hyperref}
\usepackage{xcolor}
\usepackage{listings}

\definecolor{unicaucaverde}{RGB}{0,90,60}
\hypersetup{colorlinks=true,linkcolor=unicaucaverde,urlcolor=unicaucaverde}

\titleformat{\section}{\normalfont\Large\bfseries\color{unicaucaverde}}{\thesection}{1em}{}
\titleformat{\subsection}{\normalfont\large\bfseries}{\thesubsection}{1em}{}

\theoremstyle{definition}
\newtheorem{definicion}{Definición}[section]
\newtheorem{ejemplo}[definicion]{Ejemplo}
\newtheorem{propiedad}[definicion]{Propiedad}

\lstdefinestyle{cstyle}{
  language=C,
  basicstyle=\ttfamily\small,
  keywordstyle=\color{unicaucaverde}\bfseries,
  commentstyle=\color{gray}\itshape,
  stringstyle=\color{orange!70!black},
  numbers=left,
  numberstyle=\tiny\color{gray},
  frame=single,
  breaklines=true,
  showstringspaces=false,
  tabsize=2,
  captionpos=b
}
\lstset{style=cstyle}

\title{\textcolor{unicaucaverde}{\Large Notas de clase --- Semana 2}\\[4pt]
  \large Estructuras de control: condicionales y ciclos}
\author{Programación --- Universidad del Cauca}
\date{2026}

\begin{document}
\maketitle
\tableofcontents

\section{Estructuras condicionales: \texttt{if}, \texttt{if-else}, \texttt{else if}, operador ternario}

\begin{definicion}[Sentencia condicional]
Una sentencia condicional ejecuta un bloque de código solo si una expresión --- la
\emph{condición} --- es verdadera. En C, cualquier valor distinto de \texttt{0} se considera
verdadero, y \texttt{0} se considera falso.
\end{definicion}

\begin{ejemplo}[\texttt{if} simple y \texttt{if-else}]
\begin{lstlisting}
int numero = -5;

if (numero >= 0) {
    printf("El numero es positivo o cero\n");
} else {
    printf("El numero es negativo\n");
}
\end{lstlisting}
Si la condición \texttt{numero >= 0} es falsa, se ejecuta el bloque \texttt{else}. El bloque
\texttt{else} es opcional.
\end{ejemplo}

\begin{ejemplo}[\texttt{else if} encadenado]
\begin{lstlisting}
float nota = 3.2f;

if (nota >= 4.5f) {
    printf("Excelente\n");
} else if (nota >= 3.0f) {
    printf("Aprobado\n");
} else {
    printf("Reprobado\n");
}
\end{lstlisting}
Las condiciones se evalúan en orden; se ejecuta el primer bloque cuya condición sea verdadera, y
se ignoran las siguientes.
\end{ejemplo}

\begin{definicion}[Operador ternario]
La expresión \texttt{condicion ? valor\_si\_verdadero : valor\_si\_falso} es una forma compacta de
un \texttt{if-else} que \emph{produce un valor}, útil para asignaciones simples.
\end{definicion}

\begin{ejemplo}[Operador ternario]
\begin{lstlisting}
int a = 8, b = 3, mayor;
mayor = (a > b) ? a : b;
printf("El mayor es %d\n", mayor);   // imprime 8
\end{lstlisting}
\end{ejemplo}

\begin{propiedad}[Condiciones compuestas]
Los operadores lógicos \texttt{\&\&} (y), \texttt{||} (o) y \texttt{!} (no) combinan varias
condiciones relacionales en una sola expresión. C evalúa \texttt{\&\&} y \texttt{||} de izquierda a
derecha y se detiene apenas el resultado queda determinado (\emph{cortocircuito}): en
\texttt{a \&\& b}, si \texttt{a} es falso, \texttt{b} ni siquiera se evalúa.
\end{propiedad}

\begin{ejemplo}[Condición compuesta]
\begin{lstlisting}
int edad = 20;
int tiene_carnet = 1;   // 1 = verdadero, 0 = falso

if (edad >= 18 && tiene_carnet) {
    printf("Puede ingresar\n");
} else {
    printf("No puede ingresar\n");
}
\end{lstlisting}
\end{ejemplo}

\begin{ejemplo}[Clasificador con varias condiciones compuestas]
Cuando hay más de dos posibles resultados, se encadenan varios \texttt{if}/\texttt{else if}, cada
uno con su propia condición compuesta, evaluados en orden de más a menos restrictivo:
\begin{lstlisting}
int tiene_sensor, tiene_actuador, conectado;
float bateria;   // porcentaje, 0.0 a 100.0

printf("Tiene sensor (1/0): ");
scanf("%d", &tiene_sensor);
printf("Tiene actuador (1/0): ");
scanf("%d", &tiene_actuador);
printf("Nivel de bateria: ");
scanf("%f", &bateria);
printf("Conectado a la red (1/0): ");
scanf("%d", &conectado);

if (tiene_sensor && tiene_actuador && bateria > 50 && conectado) {
    printf("Dispositivo IoT avanzado\n");
} else if (tiene_sensor && bateria >= 20 && conectado) {
    printf("Dispositivo IoT basico\n");
} else if ((tiene_sensor || tiene_actuador) && !conectado) {
    printf("Dispositivo sin conexion\n");
} else if (bateria < 10) {
    printf("Bateria critica\n");
} else {
    printf("Dispositivo no clasificado\n");
}
\end{lstlisting}
Cada rama se evalúa solo si todas las anteriores fueron falsas; por eso el orden importa: la
condición más específica (avanzado) va primero, y la más general (\texttt{else} final) queda de
respaldo para cualquier combinación no prevista.
\end{ejemplo}

\section{La instrucción \texttt{switch-case}}

\begin{definicion}[\texttt{switch-case}]
\texttt{switch} compara el valor de una expresión entera (o \texttt{char}) contra una lista de
casos (\texttt{case}), y ejecuta el bloque del primer caso que coincide. Sin un \texttt{break} al
final de cada caso, la ejecución \emph{cae} al siguiente caso (\emph{fall-through}). \texttt{default}
es opcional y se ejecuta si ningún caso coincide.
\end{definicion}

\begin{ejemplo}[Menú con \texttt{switch}]
\begin{lstlisting}
int opcion = 2;
float a = 10, b = 4, resultado;

switch (opcion) {
    case 1:
        resultado = a + b;
        printf("Suma = %.2f\n", resultado);
        break;
    case 2:
        resultado = a - b;
        printf("Resta = %.2f\n", resultado);
        break;
    case 3:
        resultado = a * b;
        printf("Producto = %.2f\n", resultado);
        break;
    default:
        printf("Opcion invalida\n");
}
\end{lstlisting}
Con \texttt{opcion = 2}, se ejecuta el caso 2 e imprime \texttt{Resta = 6.00}; el \texttt{break}
evita que también se ejecute el caso 3.
\end{ejemplo}

\section{Ciclos: \texttt{while}, \texttt{do-while}, \texttt{for}}

\begin{definicion}[Ciclo o bucle]
Un ciclo repite un bloque de instrucciones mientras se cumpla una condición. C ofrece tres formas,
que son intercambiables pero cada una es más natural según el problema.
\end{definicion}

\begin{propiedad}[\texttt{while} --- condición al inicio]
\begin{lstlisting}
int contador = 1;
while (contador <= 5) {
    printf("%d\n", contador);
    contador++;
}
\end{lstlisting}
La condición se evalúa \emph{antes} de cada repetición; si es falsa desde el principio, el cuerpo
no se ejecuta ni una vez. Útil cuando no se sabe de antemano cuántas repeticiones habrá (por
ejemplo, leer datos hasta que el usuario ingrese un centinela).
\end{propiedad}

\begin{propiedad}[\texttt{do-while} --- condición al final]
\begin{lstlisting}
int opcion;
do {
    printf("Ingrese una opcion (1-3): ");
    scanf("%d", &opcion);
} while (opcion < 1 || opcion > 3);
\end{lstlisting}
El cuerpo se ejecuta \emph{al menos una vez}, porque la condición se evalúa \emph{después}. Es la
estructura natural para validar una entrada: hay que pedirla al menos una vez antes de poder
validarla.
\end{propiedad}

\begin{propiedad}[\texttt{for} --- conteo explícito]
\begin{lstlisting}
for (int i = 1; i <= 5; i++) {
    printf("%d\n", i);
}
\end{lstlisting}
El \texttt{for} agrupa en un solo lugar la inicialización (\texttt{i = 1}), la condición
(\texttt{i <= 5}) y el incremento (\texttt{i++}). Es la forma más natural cuando se conoce de
antemano el número de repeticiones (o un rango claro).
\end{propiedad}

\section{\texttt{break}, \texttt{continue} y anidamiento}

\begin{definicion}[\texttt{break} y \texttt{continue}]
\texttt{break} termina inmediatamente el ciclo (o el \texttt{switch}) que lo contiene.
\texttt{continue} salta el resto del cuerpo del ciclo actual y pasa directamente a evaluar la
siguiente repetición.
\end{definicion}

\begin{ejemplo}[\texttt{break} y \texttt{continue}]
\begin{lstlisting}
for (int i = 1; i <= 10; i++) {
    if (i == 7) {
        break;              // termina el ciclo al llegar a 7
    }
    if (i % 2 == 0) {
        continue;           // salta los pares, no los imprime
    }
    printf("%d\n", i);      // imprime 1, 3, 5
}
\end{lstlisting}
\end{ejemplo}

\begin{definicion}[Ciclos anidados]
Un ciclo anidado es un ciclo dentro de otro. Por cada repetición del ciclo externo, el ciclo
interno se ejecuta completo. Se usan, por ejemplo, para recorrer una tabla de dos dimensiones (fila
por fila, y dentro de cada fila, columna por columna).
\end{definicion}

\begin{ejemplo}[Tabla de multiplicar con ciclos anidados]
\begin{lstlisting}
for (int fila = 1; fila <= 5; fila++) {
    for (int col = 1; col <= 5; col++) {
        printf("%3d", fila * col);
    }
    printf("\n");
}
\end{lstlisting}
Por cada valor de \texttt{fila} (ciclo externo), el ciclo interno recorre \texttt{col} de 1 a 5 e
imprime el producto; al terminar el ciclo interno, se imprime un salto de línea para pasar a la
siguiente fila.
\end{ejemplo}

\section{Matrices: llenado e indexación con ciclos anidados (adelanto)}

\begin{definicion}[Matriz (arreglo bidimensional)]
Esto es un adelanto informal: el tema completo de arreglos se estudia en una unidad posterior. Aquí
se introduce solo lo necesario para practicar ciclos anidados sobre una estructura de dos
dimensiones. Una \emph{matriz} en C se declara indicando su tipo y su número de filas y columnas:
\begin{lstlisting}
int matriz[FILAS][COLUMNAS];
\end{lstlisting}
Cada elemento se accede e identifica con dos índices, \texttt{matriz[fila][columna]}, ambos
empezando en 0: el primer elemento es \texttt{matriz[0][0]}, y el último es
\texttt{matriz[FILAS-1][COLUMNAS-1]}.
\end{definicion}

\begin{ejemplo}[Llenar una matriz con el producto de sus índices]
\begin{lstlisting}
#include <stdio.h>

int main(void) {
    int matriz[4][4];

    // Llenado: un ciclo anidado ASIGNA un valor a cada posicion
    for (int fila = 0; fila < 4; fila++) {
        for (int col = 0; col < 4; col++) {
            matriz[fila][col] = (fila + 1) * (col + 1);
        }
    }

    // Uso/impresion: un SEGUNDO ciclo anidado LEE lo que ya se guardo
    for (int fila = 0; fila < 4; fila++) {
        for (int col = 0; col < 4; col++) {
            printf("%4d", matriz[fila][col]);
        }
        printf("\n");
    }
    return 0;
}
\end{lstlisting}
El primer par de ciclos \emph{llena} la matriz (asigna un valor a cada \texttt{matriz[fila][col]});
el segundo par la \emph{recorre} para imprimirla. Llenar y usar una matriz son dos pasos distintos,
aunque el patrón de los ciclos (fila por fila, columna por columna) sea igual en ambos.
\end{ejemplo}

\begin{ejemplo}[Llenar una matriz identidad comparando índices]
\begin{lstlisting}
#include <stdio.h>

int main(void) {
    int identidad[5][5];

    for (int fila = 0; fila < 5; fila++) {
        for (int col = 0; col < 5; col++) {
            if (fila == col) {
                identidad[fila][col] = 1;   // diagonal principal
            } else {
                identidad[fila][col] = 0;
            }
        }
    }

    for (int fila = 0; fila < 5; fila++) {
        for (int col = 0; col < 5; col++) {
            printf("%2d", identidad[fila][col]);
        }
        printf("\n");
    }
    return 0;
}
\end{lstlisting}
Aquí el \emph{índice} de cada posición decide qué valor se guarda ahí: cuando \texttt{fila} y
\texttt{col} coinciden (la diagonal principal), se guarda 1; en cualquier otra posición, 0.
\end{ejemplo}

\begin{ejemplo}[Llenar una matriz leída por teclado y acumular su suma]
\begin{lstlisting}
#include <stdio.h>

int main(void) {
    int datos[3][3];
    int suma = 0;

    for (int fila = 0; fila < 3; fila++) {
        for (int col = 0; col < 3; col++) {
            printf("Valor [%d][%d]: ", fila, col);
            scanf("%d", &datos[fila][col]);
            suma += datos[fila][col];   // se acumula MIENTRAS se llena
        }
    }
    printf("Suma total = %d\n", suma);
    return 0;
}
\end{lstlisting}
No hace falta un tercer ciclo para sumar: como el acumulador \texttt{suma} vive fuera de los
ciclos, se puede actualizar dentro del mismo ciclo que llena la matriz, en el mismo momento en que
cada valor queda disponible.
\end{ejemplo}

\begin{propiedad}[El orden de los ciclos decide el orden de recorrido, no el de almacenamiento]
Una vez llena, la matriz guarda sus valores en posiciones fijas (\texttt{matriz[fila][col]}); pero
el \emph{orden en que se visita} cada posición depende de cuál ciclo es el externo y cuál el
interno. Intercambiar los ciclos externo/interno --- recorrer primero por columnas y, dentro de
cada columna, por filas --- visita las mismas posiciones pero en otro orden:
\begin{lstlisting}
// Recorrido por columnas (col es el ciclo externo)
for (int col = 0; col < 4; col++) {
    for (int fila = 0; fila < 4; fila++) {
        printf("%4d", matriz[fila][col]);
    }
    printf("\n");
}
\end{lstlisting}
Los valores impresos son los mismos que guarda la matriz; lo que cambia es el orden en que
aparecen en la salida.
\end{propiedad}

\section{Ejercicios propuestos}

\begin{enumerate}[leftmargin=1.2cm]
  \item Escribe un programa que lea un número entero e imprima si es positivo, negativo o cero.
  \item Escribe un programa que lea una edad e imprima, con \texttt{if-else if}, si la persona es
        menor de edad ($< 18$), adulto (18--59) o adulto mayor ($\geq 60$).
  \item Escribe una calculadora básica: lee dos números y un operador (\texttt{'+'}, \texttt{'-'},
        \texttt{'*'} o \texttt{'/'}) como \texttt{char}, y con \texttt{switch} imprime el resultado
        de la operación correspondiente.
  \item Escribe un programa que imprima los números del 1 al 20: primero usando \texttt{for}, y
        luego, en un segundo bloque, usando \texttt{while}.
  \item Escribe un programa que lea un entero \texttt{n} y calcule su factorial ($n!$) usando un
        ciclo.
  \item Escribe un programa que lea un entero mayor que 1 y determine, usando un ciclo, si es
        primo.
  \item Escribe un programa que imprima la tabla de multiplicar del 1 al 10 usando dos ciclos
        \texttt{for} anidados (una fila por número, del 1 al 10).
  \item Escribe un programa que, usando \texttt{do-while}, lea números enteros y los sume hasta que
        el usuario ingrese un $0$ (el $0$ no se suma), e imprima el total al final.
  \item Declara una matriz \texttt{int} de $4\times4$ y llénala, con ciclos anidados, con la suma
        de sus índices (\texttt{fila + col}); luego imprímela con otro par de ciclos anidados.
  \item Declara una matriz \texttt{int} de $n\times n$ (con $n$ leído por teclado, hasta 10) y
        llénala con una matriz identidad: 1 en la diagonal principal, 0 en el resto.
  \item Declara una matriz \texttt{float} de $3\times3$, llénala leyendo sus 9 valores por teclado,
        y calcula e imprime la suma de cada fila por separado (un total por fila, no un único total
        general).
  \item Dada una matriz \texttt{int} de $4\times4$ ya llena con el producto de sus índices (como en
        el primer ejemplo de \S 5), recórrela e imprímela \textbf{por columnas} (el ciclo externo
        avanza en \texttt{col}, el interno en \texttt{fila}), y compara el resultado impreso con el
        de recorrerla por filas.
\end{enumerate}

\section{Soluciones de los ejercicios propuestos}

\begin{description}

\item[\textbf{Ejercicio 1.}]
\begin{lstlisting}
#include <stdio.h>

int main(void) {
    int n;
    printf("Ingrese un numero: ");
    scanf("%d", &n);

    if (n > 0) {
        printf("Es positivo\n");
    } else if (n < 0) {
        printf("Es negativo\n");
    } else {
        printf("Es cero\n");
    }
    return 0;
}
\end{lstlisting}

\item[\textbf{Ejercicio 2.}]
\begin{lstlisting}
#include <stdio.h>

int main(void) {
    int edad;
    printf("Ingrese la edad: ");
    scanf("%d", &edad);

    if (edad < 18) {
        printf("Menor de edad\n");
    } else if (edad < 60) {
        printf("Adulto\n");
    } else {
        printf("Adulto mayor\n");
    }
    return 0;
}
\end{lstlisting}

\item[\textbf{Ejercicio 3.}]
\begin{lstlisting}
#include <stdio.h>

int main(void) {
    float a, b, resultado;
    char operador;

    printf("Ingrese numero, operador, numero (ej: 4 + 2): ");
    scanf("%f %c %f", &a, &operador, &b);

    switch (operador) {
        case '+':
            resultado = a + b;
            break;
        case '-':
            resultado = a - b;
            break;
        case '*':
            resultado = a * b;
            break;
        case '/':
            if (b == 0) {
                printf("Error: division por cero\n");
                return 1;
            }
            resultado = a / b;
            break;
        default:
            printf("Operador invalido\n");
            return 1;
    }
    printf("Resultado = %.2f\n", resultado);
    return 0;
}
\end{lstlisting}

\item[\textbf{Ejercicio 4.}]
\begin{lstlisting}
#include <stdio.h>

int main(void) {
    printf("Con for:\n");
    for (int i = 1; i <= 20; i++) {
        printf("%d ", i);
    }
    printf("\n");

    printf("Con while:\n");
    int i = 1;
    while (i <= 20) {
        printf("%d ", i);
        i++;
    }
    printf("\n");
    return 0;
}
\end{lstlisting}

\item[\textbf{Ejercicio 5.}]
\begin{lstlisting}
#include <stdio.h>

int main(void) {
    int n;
    long factorial = 1;

    printf("Ingrese n: ");
    scanf("%d", &n);

    for (int i = 1; i <= n; i++) {
        factorial *= i;
    }
    printf("%d! = %ld\n", n, factorial);
    return 0;
}
\end{lstlisting}
Con $n=5$: $factorial = 1 \cdot 2 \cdot 3 \cdot 4 \cdot 5 = 120$.

\item[\textbf{Ejercicio 6.}]
\begin{lstlisting}
#include <stdio.h>

int main(void) {
    int n, es_primo = 1;

    printf("Ingrese un numero mayor que 1: ");
    scanf("%d", &n);

    for (int i = 2; i < n; i++) {
        if (n % i == 0) {
            es_primo = 0;
            break;          // ya encontramos un divisor, no hace falta seguir
        }
    }

    if (es_primo) {
        printf("%d es primo\n", n);
    } else {
        printf("%d no es primo\n", n);
    }
    return 0;
}
\end{lstlisting}
En cuanto se encuentra un divisor de \texttt{n} distinto de 1 y de sí mismo, \texttt{break} corta
el ciclo: no hace falta seguir probando.

\item[\textbf{Ejercicio 7.}]
\begin{lstlisting}
#include <stdio.h>

int main(void) {
    for (int fila = 1; fila <= 10; fila++) {
        for (int col = 1; col <= 10; col++) {
            printf("%4d", fila * col);
        }
        printf("\n");
    }
    return 0;
}
\end{lstlisting}

\item[\textbf{Ejercicio 8.}]
\begin{lstlisting}
#include <stdio.h>

int main(void) {
    int numero, suma = 0;

    do {
        printf("Ingrese un numero (0 para terminar): ");
        scanf("%d", &numero);
        if (numero != 0) {
            suma += numero;
        }
    } while (numero != 0);

    printf("Suma total = %d\n", suma);
    return 0;
}
\end{lstlisting}
Se usa \texttt{do-while} porque hay que pedir el primer número antes de poder evaluar si es el
centinela (\texttt{0}) que termina el ciclo.

\item[\textbf{Ejercicio 9.}]
\begin{lstlisting}
#include <stdio.h>

int main(void) {
    int matriz[4][4];

    for (int fila = 0; fila < 4; fila++) {
        for (int col = 0; col < 4; col++) {
            matriz[fila][col] = fila + col;
        }
    }

    for (int fila = 0; fila < 4; fila++) {
        for (int col = 0; col < 4; col++) {
            printf("%3d", matriz[fila][col]);
        }
        printf("\n");
    }
    return 0;
}
\end{lstlisting}

\item[\textbf{Ejercicio 10.}]
\begin{lstlisting}
#include <stdio.h>

int main(void) {
    int n;
    int identidad[10][10];

    printf("Tamano n (hasta 10): ");
    scanf("%d", &n);

    for (int fila = 0; fila < n; fila++) {
        for (int col = 0; col < n; col++) {
            identidad[fila][col] = (fila == col) ? 1 : 0;
        }
    }

    for (int fila = 0; fila < n; fila++) {
        for (int col = 0; col < n; col++) {
            printf("%2d", identidad[fila][col]);
        }
        printf("\n");
    }
    return 0;
}
\end{lstlisting}

\item[\textbf{Ejercicio 11.}]
\begin{lstlisting}
#include <stdio.h>

int main(void) {
    float valores[3][3];

    for (int fila = 0; fila < 3; fila++) {
        for (int col = 0; col < 3; col++) {
            printf("Valor [%d][%d]: ", fila, col);
            scanf("%f", &valores[fila][col]);
        }
    }

    for (int fila = 0; fila < 3; fila++) {
        float sumaFila = 0;
        for (int col = 0; col < 3; col++) {
            sumaFila += valores[fila][col];
        }
        printf("Suma fila %d = %.2f\n", fila, sumaFila);
    }
    return 0;
}
\end{lstlisting}
La suma de cada fila se reinicia (\texttt{sumaFila = 0}) antes de cada nueva fila: son 3 sumas
independientes, no una acumulada entre filas.

\item[\textbf{Ejercicio 12.}]
\begin{lstlisting}
#include <stdio.h>

int main(void) {
    int matriz[4][4];

    for (int fila = 0; fila < 4; fila++) {
        for (int col = 0; col < 4; col++) {
            matriz[fila][col] = (fila + 1) * (col + 1);
        }
    }

    printf("Por filas:\n");
    for (int fila = 0; fila < 4; fila++) {
        for (int col = 0; col < 4; col++) {
            printf("%4d", matriz[fila][col]);
        }
        printf("\n");
    }

    printf("Por columnas:\n");
    for (int col = 0; col < 4; col++) {
        for (int fila = 0; fila < 4; fila++) {
            printf("%4d", matriz[fila][col]);
        }
        printf("\n");
    }
    return 0;
}
\end{lstlisting}
La matriz almacenada es la misma en ambos casos; lo único que cambia es el orden de recorrido:
"por filas" imprime primero toda la fila 0, luego toda la fila 1, etc.; "por columnas" imprime
primero toda la columna 0 (un valor de cada fila), luego toda la columna 1, etc. Como en este
ejemplo \texttt{matriz[fila][col] = (fila+1)*(col+1)} es simétrica respecto a fila y columna, las
dos salidas se ven parecidas; con una matriz no simétrica (como la del Ejercicio 9) el cambio de
orden sí produce una salida visiblemente distinta.

\end{description}

\end{document}
